\documentclass[a4paper,12pt]{article}
\usepackage[italian]{babel}
\usepackage[utf8]{inputenc}
\usepackage[T1]{fontenc}
\usepackage{geometry}
\usepackage{setspace}
\geometry{margin=2.5cm}
\onehalfspacing

\title{Soluzioni - Verifica scritta di Informatica - SQL e Algebra Relazionale}
\author{Prof. Fedeli Massimo}
\date{}

\begin{document}
	\maketitle
	
	\section*{Schema logico}
	
	\begin{flushleft}
		UTENTI(id\_utente PK, cognome, nome, classe, email, data\_iscrizione) \\
		LIBRI(id\_libro PK, titolo, autore, genere, anno\_pubblicazione, editore, copie\_totali) \\
		PRESTITI(id\_prestito PK, id\_utente FK, id\_libro FK, data\_prestito, data\_scadenza, data\_restituzione) \\
		PRENOTAZIONI(id\_prenotazione PK, id\_utente FK, id\_libro FK, data\_prenotazione, stato)
	\end{flushleft}
	
	\section*{Soluzioni}
	
	\begin{enumerate}
		\item
\begin{verbatim}
SELECT cognome, nome, classe
FROM UTENTI
ORDER BY classe, cognome, nome;
\end{verbatim}

		\item
\begin{verbatim}
SELECT titolo, autore, genere
FROM LIBRI
WHERE anno_pubblicazione > 2018;
\end{verbatim}

		\item
\begin{verbatim}
SELECT u.nome, u.cognome, l.titolo, p.data_scadenza
FROM PRESTITI p
JOIN UTENTI u ON p.id_utente = u.id_utente
JOIN LIBRI l ON p.id_libro = l.id_libro
WHERE p.data_restituzione IS NULL;
\end{verbatim}

		\item
\begin{verbatim}
SELECT titolo, autore, genere
FROM LIBRI
WHERE genere IN ('Informatica', 'Matematica')
ORDER BY autore, titolo;
\end{verbatim}

		\item
\begin{verbatim}
SELECT genere, COUNT(*) AS numero_libri
FROM LIBRI
GROUP BY genere;
\end{verbatim}

		\item
\begin{verbatim}
SELECT u.id_utente, u.cognome, u.nome, COUNT(p.id_prestito) AS totale_prestiti
FROM UTENTI u
LEFT JOIN PRESTITI p ON u.id_utente = p.id_utente
GROUP BY u.id_utente, u.cognome, u.nome;
\end{verbatim}

		\item
\begin{verbatim}
SELECT l.titolo
FROM LIBRI l
LEFT JOIN PRESTITI p ON l.id_libro = p.id_libro
WHERE p.id_prestito IS NULL;
\end{verbatim}

		\item
\begin{verbatim}
SELECT u.id_utente, u.cognome, u.nome
FROM UTENTI u
JOIN PRESTITI p ON u.id_utente = p.id_utente
GROUP BY u.id_utente, u.cognome, u.nome
HAVING COUNT(DISTINCT p.id_libro) >= 3;
\end{verbatim}

		\item
\begin{verbatim}
SELECT l.id_libro, l.titolo, COUNT(*) AS prenotazioni_attive
FROM LIBRI l
JOIN PRENOTAZIONI pr ON l.id_libro = pr.id_libro
WHERE pr.stato = 'attiva'
GROUP BY l.id_libro, l.titolo
HAVING COUNT(*) > 2;
\end{verbatim}

		\item
\begin{verbatim}
SELECT u.id_utente, u.cognome, u.nome, COUNT(*) AS totale_prestiti
FROM UTENTI u
JOIN PRESTITI p ON u.id_utente = p.id_utente
GROUP BY u.id_utente, u.cognome, u.nome
HAVING COUNT(*) = (
    SELECT MAX(tot)
    FROM (
        SELECT COUNT(*) AS tot
        FROM PRESTITI
        GROUP BY id_utente
    ) AS conteggi
);
\end{verbatim}

		\item
\begin{verbatim}
SELECT l.id_libro, l.titolo, l.copie_totali,
       COUNT(p.id_prestito) AS copie_attualmente_in_prestito
FROM LIBRI l
LEFT JOIN PRESTITI p
  ON l.id_libro = p.id_libro
 AND p.data_restituzione IS NULL
GROUP BY l.id_libro, l.titolo, l.copie_totali
HAVING COUNT(p.id_prestito) < l.copie_totali;
\end{verbatim}

		\item
\begin{verbatim}
PI cognome, nome (
  UTENTI JOIN_{UTENTI.id_utente = PRESTITI.id_utente}
  SIGMA data_restituzione IS NULL (PRESTITI)
)
\end{verbatim}

		\item
\begin{verbatim}
PI titolo (
  LIBRI JOIN_{LIBRI.id_libro = PRENOTAZIONI.id_libro} PRENOTAZIONI
)
\end{verbatim}

		\item
\begin{verbatim}
SELECT u.id_utente, u.cognome, u.nome
FROM UTENTI u
WHERE u.id_utente IN (
    SELECT id_utente
    FROM PRESTITI
)
AND u.id_utente NOT IN (
    SELECT id_utente
    FROM PRENOTAZIONI
);
\end{verbatim}

		\item
\begin{verbatim}
SELECT genere, AVG(anno_pubblicazione) AS anno_medio
FROM LIBRI
GROUP BY genere;
\end{verbatim}
	\end{enumerate}
	
\end{document}
